\documentclass[11pt]{article}

\usepackage[margin=1in]{geometry}
\usepackage{enumitem}
\usepackage{amsmath, amssymb}

\setlist[enumerate]{itemsep=0.6em, topsep=0.6em}

\begin{document}

\begin{center}
  {\Large \textbf{CS 426 Examination}}\\[0.25em]
  {\large \textbf{Large-Scale System Design and Architecture -- Answer Key}}\\[0.5em]
\end{center}

\begin{enumerate}[label=\textbf{\arabic*.}]

  \item \textbf{(B)} Write-through caching.\\
  In write-through caching, every write operation updates both cache and database synchronously. 
  This ensures cache and database are always consistent, though it adds latency to write operations. 
  Write-behind is faster but risks data loss; cache-aside can have stale data.

  \item \textbf{(C)} Any two of the three properties, but not all three.\\
  The CAP theorem states that in the presence of network partitions (P), you must choose between 
  consistency (C) and availability (A). Most systems choose CP (consistent but may be unavailable) 
  or AP (available but eventually consistent). CA is theoretically possible but unrealistic since 
  network partitions are inevitable.

  \item \textbf{(B)} To distribute incoming requests across multiple servers for scalability and 
  reliability.\\
  Load balancers distribute traffic using algorithms like round-robin, least connections, or 
  weighted distribution. This prevents any single server from being overwhelmed and provides 
  fault tolerance—if one server fails, traffic routes to healthy servers.

  \item \textbf{(B)} Hash-based sharding.\\
  Hash-based sharding applies a hash function to the shard key (e.g., user ID) to determine which 
  shard stores the data: $\text{shard} = \text{hash}(\text{key}) \mod N$. This distributes data 
  relatively evenly but makes range queries difficult.

  \item \emph{Sample answer:}\\
  \textbf{Eventual consistency}: After an update, all replicas will eventually converge to the 
  same value, but they may temporarily differ. Reads might return stale data.

  \textbf{Trade-offs}:
  \begin{itemize}
    \item \textbf{Strong consistency}: All reads see the latest write. Requires coordination (locks, 
    consensus), reducing availability and increasing latency
    \item \textbf{Eventual consistency}: Higher availability and performance, but complexity in 
    handling conflicts and stale reads
  \end{itemize}

  \textbf{Choose eventual consistency when}:
  \begin{itemize}
    \item Availability is critical (e.g., social media feeds, shopping carts)
    \item Temporary staleness is acceptable
    \item System must tolerate network partitions
    \item Examples: DynamoDB, Cassandra, DNS
  \end{itemize}

  \item \emph{Sample answer:}\\
  \textbf{Microservices}: Architecture where application is composed of small, independently 
  deployable services, each responsible for a specific business capability.

  \textbf{Advantages over monolith}:
  \begin{itemize}
    \item Independent scaling of services based on demand
    \item Technology diversity (each service can use different languages/frameworks)
    \item Faster deployment cycles (teams can release independently)
    \item Fault isolation (failure in one service doesn't crash entire system)
  \end{itemize}

  \textbf{Additional complexities}:
  \begin{itemize}
    \item \textbf{Deployment}: Need orchestration (Kubernetes), service discovery, configuration management
    \item \textbf{Debugging}: Distributed tracing required, harder to track requests across services
    \item \textbf{Communication}: Network calls are slower and less reliable than in-process calls; 
    need to handle failures, timeouts, retries
    \item \textbf{Data consistency}: Managing transactions across services is challenging
  \end{itemize}

  \item \emph{Sample answer:}\\
  \textbf{Distributed message queue}: Intermediary that decouples producers (who send messages) 
  from consumers (who process them), storing messages until consumed.

  \textbf{Purpose and benefits}:
  \begin{itemize}
    \item \textbf{Asynchronous processing}: Producers don't wait for consumers, improving responsiveness
    \item \textbf{Load leveling}: Queue absorbs traffic spikes; consumers process at their own pace
    \item \textbf{Reliability}: Messages persist in queue even if consumers fail
    \item \textbf{Scalability}: Multiple consumers can process messages in parallel
  \end{itemize}

  \textbf{Implementation (Kafka example)}:
  \begin{itemize}
    \item Messages organized into topics (categories)
    \item Topics divided into partitions for parallelism
    \item Consumer groups allow multiple consumers to process messages
    \item Provides ordering guarantees within partitions
  \end{itemize}

  \textbf{Patterns enabled}:
  \begin{itemize}
    \item Event-driven architecture, pub-sub messaging, task queues, stream processing
  \end{itemize}

  \item \emph{Sample answer:}\\
  \textbf{Token bucket algorithm}:
  \begin{itemize}
    \item Bucket holds tokens; tokens added at constant rate
    \item Each request consumes one token; rejected if bucket empty
    \item Allows bursts up to bucket capacity
    \item Good for APIs with occasional traffic spikes
  \end{itemize}

  \textbf{Leaky bucket algorithm}:
  \begin{itemize}
    \item Requests enter bucket; processed at constant rate (``leak'')
    \item If bucket full, new requests rejected
    \item Smooths traffic, prevents bursts
    \item Better for strict rate enforcement
  \end{itemize}

  \textbf{Distributed implementation}:
  \begin{itemize}
    \item \textbf{Centralized counter}: Use Redis with atomic operations (INCR, EXPIRE) to track 
    request counts per user/IP. Fast but single point of failure
    \item \textbf{Distributed coordination}: Use consistent hashing to route user requests to 
    specific rate-limiting servers; reduces contention
    \item \textbf{Approximate counting}: Each server tracks counts locally; periodically sync. 
    Faster but less accurate
    \item \textbf{Sliding window}: Use sorted sets in Redis to track timestamps of recent requests
  \end{itemize}

  Most production systems use Redis-based token bucket with sliding windows for accuracy and performance.

\end{enumerate}

\end{document}
